\chapter*{\center \Large ANEXOS}
\addcontentsline{toc}{section}{\bfseries ANEXOS}
\markboth{ANEXOS}{ANEXOS}

\section*{Anexo A. Tablas completas por episodio}

Las tablas siguientes listan, para cada escenario, los parámetros de
entrada y las métricas de salida de todos los episodios de
entrenamiento. Se generan directamente desde
\texttt{resultados\_<escenario>.csv} con el script
\texttt{tablas\_tex.py}, de modo que cada número es reproducible con
los comandos del plan de pruebas. Q ini.\ es el número de estados
distintos que la tabla Q ya conocía al empezar el episodio; $r$/dec.\ es
la recompensa media por decisión, medida cada 5~s con el mismo reloj en
todos los controles.

% Generado por tablas_tex.py a partir de los CSV del proyecto. No editar a mano.
\begin{table}[H]
\setstretch{1}
\centering
\scriptsize
\setlength{\tabcolsep}{3pt}
\begin{tabular}{rrrrrrrrrrrr}
\toprule
\multicolumn{4}{c}{Entradas del episodio} & \multicolumn{6}{c}{Salidas del episodio (entrenamiento, $\varepsilon>0$)} & \multicolumn{2}{c}{Política congelada} \\
\cmidrule(lr){1-4}\cmidrule(lr){5-10}\cmidrule(lr){11-12}
Ep. & $\varepsilon$ & Semilla & Q ini. & Retraso & Espera & Cola & $r$/dec. & Cambios & Paradas & Retraso & Espera \\
 &  &  & (estados) & (s/veh) & (s/veh) & (veh) &  & de fase & (/veh) & (s/veh) & (s/veh) \\
\midrule
fijo & 0 & 42 & -- & 22.83 & 9.30 & 5.16 & $-$5.58 & 131 & 0.67 & -- & -- \\
\midrule
1 & 1.000 & 43 & 0 & 80.14 & 30.84 & 14.82 & $-$15.98 & 227 & 3.16 & -- & -- \\
2 & 0.900 & 44 & 81 & 51.59 & 21.21 & 10.90 & $-$11.64 & 219 & 2.16 & -- & -- \\
3 & 0.810 & 45 & 104 & 36.22 & 15.46 & 7.97 & $-$8.51 & 215 & 1.43 & -- & -- \\
4 & 0.729 & 46 & 114 & 33.46 & 14.07 & 7.51 & $-$8.00 & 209 & 1.29 & -- & -- \\
5 & 0.656 & 47 & 119 & 31.21 & 12.33 & 6.62 & $-$7.00 & 222 & 1.27 & 22.15 & 7.26 \\
6 & 0.590 & 48 & 120 & 35.50 & 13.32 & 7.19 & $-$7.61 & 223 & 1.46 & -- & -- \\
7 & 0.531 & 49 & 120 & 24.22 & 8.92 & 4.96 & $-$5.25 & 220 & 0.92 & -- & -- \\
8 & 0.478 & 50 & 123 & 29.17 & 11.04 & 6.03 & $-$6.35 & 219 & 1.17 & -- & -- \\
9 & 0.430 & 51 & 123 & 33.91 & 13.05 & 6.91 & $-$7.29 & 225 & 1.41 & -- & -- \\
10 & 0.387 & 52 & 123 & 27.00 & 9.95 & 5.45 & $-$5.72 & 222 & 1.09 & 21.43 & 7.43 \\
11 & 0.349 & 53 & 123 & 24.91 & 9.03 & 5.04 & $-$5.30 & 215 & 0.98 & -- & -- \\
12 & 0.314 & 54 & 124 & 22.93 & 8.10 & 4.56 & $-$4.74 & 230 & 0.92 & -- & -- \\
13 & 0.282 & 55 & 124 & 21.90 & 7.37 & 4.20 & $-$4.39 & 231 & 0.85 & -- & -- \\
14 & 0.254 & 56 & 124 & 22.11 & 7.51 & 4.24 & $-$4.44 & 222 & 0.86 & -- & -- \\
15 & 0.229 & 57 & 124 & 22.18 & 7.45 & 4.23 & $-$4.41 & 227 & 0.89 & 20.13 & 6.28 \\
16 & 0.206 & 58 & 124 & 20.03 & 6.40 & 3.71 & $-$3.88 & 216 & 0.74 & -- & -- \\
17 & 0.185 & 59 & 125 & 22.28 & 7.77 & 4.40 & $-$4.63 & 215 & 0.88 & -- & -- \\
18 & 0.167 & 60 & 125 & 21.78 & 7.36 & 4.20 & $-$4.39 & 223 & 0.85 & -- & -- \\
19 & 0.150 & 61 & 125 & 26.04 & 9.15 & 5.08 & $-$5.31 & 222 & 1.07 & -- & -- \\
20 & 0.135 & 62 & 126 & 22.10 & 7.59 & 4.31 & $-$4.48 & 229 & 0.88 & 39.18 & 12.66 \\
21 & 0.122 & 63 & 126 & 19.60 & 6.07 & 3.55 & $-$3.69 & 223 & 0.74 & -- & -- \\
22 & 0.109 & 64 & 126 & 21.73 & 7.21 & 4.11 & $-$4.26 & 225 & 0.86 & -- & -- \\
23 & 0.098 & 65 & 127 & 20.76 & 6.71 & 3.86 & $-$3.99 & 225 & 0.80 & -- & -- \\
24 & 0.089 & 66 & 127 & 22.09 & 7.42 & 4.21 & $-$4.37 & 219 & 0.85 & -- & -- \\
25 & 0.080 & 67 & 127 & 19.42 & 6.24 & 3.62 & $-$3.75 & 213 & 0.72 & 20.50 & 6.56 \\
26 & 0.072 & 68 & 127 & 27.84 & 9.91 & 5.45 & $-$5.69 & 230 & 1.18 & -- & -- \\
27 & 0.065 & 69 & 127 & 24.92 & 8.57 & 4.79 & $-$5.03 & 223 & 1.00 & -- & -- \\
28 & 0.058 & 70 & 128 & 21.84 & 7.22 & 4.12 & $-$4.29 & 237 & 0.86 & -- & -- \\
29 & 0.052 & 71 & 128 & 25.07 & 8.80 & 4.92 & $-$5.15 & 229 & 1.03 & -- & -- \\
30 & 0.050 & 72 & 129 & 17.80 & 5.31 & 3.16 & $-$3.30 & 215 & 0.64 & 19.62 & 6.18 \\
31 & 0.050 & 73 & 129 & 20.97 & 6.86 & 3.93 & $-$4.12 & 226 & 0.81 & -- & -- \\
32 & 0.050 & 74 & 129 & 20.66 & 6.73 & 3.83 & $-$3.98 & 229 & 0.81 & -- & -- \\
33 & 0.050 & 75 & 129 & 20.62 & 6.73 & 3.87 & $-$4.02 & 223 & 0.82 & -- & -- \\
34 & 0.050 & 76 & 130 & 19.42 & 6.02 & 3.52 & $-$3.67 & 223 & 0.74 & -- & -- \\
35 & 0.050 & 77 & 130 & 19.66 & 6.23 & 3.62 & $-$3.76 & 219 & 0.74 & 20.02 & 6.36 \\
36 & 0.050 & 78 & 130 & 24.66 & 8.57 & 4.80 & $-$5.02 & 227 & 1.03 & -- & -- \\
37 & 0.050 & 79 & 130 & 19.62 & 6.17 & 3.60 & $-$3.73 & 217 & 0.74 & -- & -- \\
38 & 0.050 & 80 & 130 & 23.06 & 7.92 & 4.47 & $-$4.64 & 225 & 0.94 & -- & -- \\
39 & 0.050 & 81 & 130 & 19.87 & 6.33 & 3.67 & $-$3.81 & 229 & 0.78 & -- & -- \\
40 & 0.050 & 82 & 130 & 21.65 & 7.13 & 4.07 & $-$4.23 & 223 & 0.84 & 21.39 & 7.05 \\
\bottomrule
\end{tabular}
\caption[Escenario cruce: episodios de entrenamiento]{Escenario cruce: parámetros de entrada y métricas de salida por episodio. 40 episodios de entrenamiento con semilla base 42; la fila fijo es el tiempo fijo con la misma semilla base. La política congelada se evalúa cada 5 episodios con una semilla fija ajena al entrenamiento y a la evaluación final.}
\label{tab:cruce-episodios}
\end{table}


% Generado por tablas_tex.py a partir de los CSV del proyecto. No editar a mano.
\begin{table}[H]
\setstretch{1}
\centering
\scriptsize
\setlength{\tabcolsep}{3pt}
\begin{tabular}{rrrrrrrrrrrr}
\toprule
\multicolumn{4}{c}{Entradas del episodio} & \multicolumn{6}{c}{Salidas del episodio (entrenamiento, $\varepsilon>0$)} & \multicolumn{2}{c}{Política congelada} \\
\cmidrule(lr){1-4}\cmidrule(lr){5-10}\cmidrule(lr){11-12}
Ep. & $\varepsilon$ & Semilla & Q ini. & Retraso & Espera & Cola & $r$/dec. & Cambios & Paradas & Retraso & Espera \\
 &  &  & (estados) & (s/veh) & (s/veh) & (veh) &  & de fase & (/veh) & (s/veh) & (s/veh) \\
\midrule
fijo & 0 & 42 & -- & 27.94 & 15.96 & 7.69 & $-$8.36 & 230 & 0.89 & -- & -- \\
\midrule
1 & 1.000 & 43 & 0 & 207.74 & 132.44 & 57.32 & $-$67.25 & 280 & 6.91 & -- & -- \\
2 & 0.960 & 44 & 106 & 207.37 & 134.95 & 57.69 & $-$67.97 & 296 & 6.30 & -- & -- \\
3 & 0.922 & 45 & 128 & 154.32 & 99.04 & 47.59 & $-$56.27 & 256 & 4.74 & -- & -- \\
4 & 0.885 & 46 & 157 & 103.72 & 67.25 & 33.75 & $-$39.39 & 248 & 2.77 & -- & -- \\
5 & 0.849 & 47 & 199 & 89.60 & 59.52 & 28.76 & $-$34.10 & 245 & 2.27 & 209.53 & 153.51 \\
6 & 0.815 & 48 & 241 & 111.92 & 70.79 & 32.27 & $-$37.34 & 284 & 3.66 & -- & -- \\
7 & 0.783 & 49 & 258 & 118.04 & 73.22 & 34.30 & $-$39.97 & 260 & 3.52 & -- & -- \\
8 & 0.751 & 50 & 270 & 59.94 & 39.21 & 18.58 & $-$21.49 & 261 & 1.66 & -- & -- \\
9 & 0.721 & 51 & 274 & 69.69 & 46.06 & 22.02 & $-$25.66 & 255 & 1.82 & -- & -- \\
10 & 0.693 & 52 & 282 & 106.36 & 67.16 & 31.43 & $-$36.27 & 280 & 3.01 & 48.71 & 30.97 \\
11 & 0.665 & 53 & 297 & 52.21 & 34.99 & 17.37 & $-$20.18 & 240 & 1.39 & -- & -- \\
12 & 0.638 & 54 & 307 & 83.47 & 53.19 & 25.56 & $-$29.38 & 267 & 2.41 & -- & -- \\
13 & 0.613 & 55 & 313 & 42.01 & 26.84 & 13.35 & $-$15.28 & 231 & 1.18 & -- & -- \\
14 & 0.588 & 56 & 320 & 47.09 & 30.37 & 15.07 & $-$17.20 & 240 & 1.31 & -- & -- \\
15 & 0.565 & 57 & 328 & 40.80 & 25.90 & 12.70 & $-$14.39 & 256 & 1.19 & 26.65 & 15.05 \\
16 & 0.542 & 58 & 340 & 39.46 & 24.78 & 12.34 & $-$13.88 & 258 & 1.18 & -- & -- \\
17 & 0.520 & 59 & 343 & 39.12 & 24.51 & 12.17 & $-$13.72 & 263 & 1.17 & -- & -- \\
18 & 0.500 & 60 & 348 & 47.08 & 29.93 & 14.80 & $-$16.72 & 263 & 1.37 & -- & -- \\
19 & 0.480 & 61 & 351 & 34.93 & 21.32 & 10.56 & $-$11.79 & 260 & 1.08 & -- & -- \\
20 & 0.460 & 62 & 355 & 38.80 & 23.88 & 11.77 & $-$13.10 & 271 & 1.20 & 33.96 & 20.13 \\
21 & 0.442 & 63 & 359 & 38.99 & 24.10 & 11.61 & $-$12.96 & 267 & 1.18 & -- & -- \\
22 & 0.424 & 64 & 361 & 31.30 & 18.44 & 8.99 & $-$9.98 & 249 & 0.98 & -- & -- \\
23 & 0.407 & 65 & 362 & 34.65 & 20.72 & 10.07 & $-$11.11 & 272 & 1.11 & -- & -- \\
24 & 0.391 & 66 & 362 & 32.28 & 19.13 & 9.39 & $-$10.35 & 259 & 1.04 & -- & -- \\
25 & 0.375 & 67 & 362 & 32.77 & 19.28 & 9.47 & $-$10.39 & 265 & 1.06 & 32.32 & 18.48 \\
26 & 0.360 & 68 & 363 & 39.31 & 24.12 & 12.02 & $-$13.35 & 281 & 1.24 & -- & -- \\
27 & 0.346 & 69 & 367 & 31.48 & 18.47 & 8.93 & $-$9.85 & 266 & 1.01 & -- & -- \\
28 & 0.332 & 70 & 368 & 29.16 & 16.68 & 8.09 & $-$8.89 & 257 & 0.97 & -- & -- \\
29 & 0.319 & 71 & 369 & 33.72 & 20.13 & 9.83 & $-$10.90 & 271 & 1.08 & -- & -- \\
30 & 0.306 & 72 & 372 & 31.39 & 18.35 & 9.02 & $-$9.90 & 266 & 1.03 & 37.14 & 21.75 \\
31 & 0.294 & 73 & 374 & 32.10 & 19.04 & 9.41 & $-$10.38 & 264 & 1.03 & -- & -- \\
32 & 0.282 & 74 & 374 & 30.87 & 17.91 & 8.72 & $-$9.59 & 272 & 1.01 & -- & -- \\
33 & 0.271 & 75 & 374 & 30.37 & 17.52 & 8.58 & $-$9.41 & 261 & 0.99 & -- & -- \\
34 & 0.260 & 76 & 374 & 29.81 & 17.27 & 8.49 & $-$9.35 & 248 & 0.96 & -- & -- \\
35 & 0.250 & 77 & 375 & 28.95 & 16.68 & 8.09 & $-$8.90 & 244 & 0.93 & 26.18 & 14.52 \\
36 & 0.240 & 78 & 375 & 29.90 & 17.22 & 8.47 & $-$9.24 & 267 & 0.99 & -- & -- \\
37 & 0.230 & 79 & 375 & 28.73 & 16.33 & 8.05 & $-$8.78 & 259 & 0.97 & -- & -- \\
38 & 0.221 & 80 & 376 & 33.45 & 19.83 & 9.48 & $-$10.41 & 279 & 1.09 & -- & -- \\
39 & 0.212 & 81 & 376 & 31.79 & 18.66 & 9.25 & $-$10.19 & 267 & 1.05 & -- & -- \\
40 & 0.204 & 82 & 378 & 27.71 & 15.55 & 7.56 & $-$8.22 & 255 & 0.93 & 26.06 & 14.63 \\
\bottomrule
\end{tabular}
\caption[Escenario cruce2: episodios de entrenamiento (parte 1 de 3)]{Escenario cruce2: parámetros de entrada y métricas de salida por episodio (parte 1 de 3). 120 episodios de entrenamiento con semilla base 42; la fila fijo es el tiempo fijo con la misma semilla base. La política congelada se evalúa cada 5 episodios con una semilla fija ajena al entrenamiento y a la evaluación final.}
\label{tab:cruce2-episodios}
\end{table}

\begin{table}[H]
\setstretch{1}
\centering
\scriptsize
\setlength{\tabcolsep}{3pt}
\begin{tabular}{rrrrrrrrrrrr}
\toprule
\multicolumn{4}{c}{Entradas del episodio} & \multicolumn{6}{c}{Salidas del episodio (entrenamiento, $\varepsilon>0$)} & \multicolumn{2}{c}{Política congelada} \\
\cmidrule(lr){1-4}\cmidrule(lr){5-10}\cmidrule(lr){11-12}
Ep. & $\varepsilon$ & Semilla & Q ini. & Retraso & Espera & Cola & $r$/dec. & Cambios & Paradas & Retraso & Espera \\
 &  &  & (estados) & (s/veh) & (s/veh) & (veh) &  & de fase & (/veh) & (s/veh) & (s/veh) \\
\midrule
41 & 0.195 & 83 & 378 & 26.85 & 15.01 & 7.37 & $-$7.99 & 248 & 0.89 & -- & -- \\
42 & 0.188 & 84 & 378 & 27.79 & 15.62 & 7.64 & $-$8.31 & 255 & 0.92 & -- & -- \\
43 & 0.180 & 85 & 378 & 30.56 & 17.63 & 8.59 & $-$9.43 & 264 & 1.00 & -- & -- \\
44 & 0.173 & 86 & 381 & 32.95 & 19.21 & 9.53 & $-$10.44 & 266 & 1.08 & -- & -- \\
45 & 0.166 & 87 & 382 & 31.11 & 18.07 & 8.84 & $-$9.68 & 271 & 1.01 & 26.33 & 14.88 \\
46 & 0.159 & 88 & 383 & 27.10 & 15.12 & 7.31 & $-$7.99 & 255 & 0.91 & -- & -- \\
47 & 0.153 & 89 & 383 & 28.02 & 15.86 & 7.57 & $-$8.27 & 247 & 0.92 & -- & -- \\
48 & 0.147 & 90 & 383 & 27.67 & 15.61 & 7.58 & $-$8.24 & 254 & 0.92 & -- & -- \\
49 & 0.141 & 91 & 383 & 28.56 & 16.27 & 7.98 & $-$8.73 & 254 & 0.94 & -- & -- \\
50 & 0.135 & 92 & 383 & 27.64 & 15.66 & 7.59 & $-$8.32 & 248 & 0.90 & 34.65 & 20.20 \\
51 & 0.130 & 93 & 383 & 26.60 & 14.81 & 7.20 & $-$7.82 & 247 & 0.88 & -- & -- \\
52 & 0.125 & 94 & 383 & 28.53 & 16.24 & 7.99 & $-$8.73 & 246 & 0.93 & -- & -- \\
53 & 0.120 & 95 & 383 & 28.74 & 16.29 & 8.04 & $-$8.77 & 252 & 0.95 & -- & -- \\
54 & 0.115 & 96 & 384 & 26.04 & 14.45 & 7.12 & $-$7.73 & 239 & 0.85 & -- & -- \\
55 & 0.110 & 97 & 384 & 27.12 & 15.23 & 7.46 & $-$8.11 & 248 & 0.91 & 32.38 & 18.73 \\
56 & 0.106 & 98 & 384 & 26.14 & 14.46 & 7.02 & $-$7.59 & 244 & 0.87 & -- & -- \\
57 & 0.102 & 99 & 385 & 26.34 & 14.69 & 7.17 & $-$7.77 & 241 & 0.88 & -- & -- \\
58 & 0.098 & 100 & 385 & 28.50 & 16.01 & 7.85 & $-$8.53 & 267 & 0.95 & -- & -- \\
59 & 0.094 & 101 & 385 & 28.46 & 16.18 & 7.93 & $-$8.69 & 251 & 0.93 & -- & -- \\
60 & 0.090 & 102 & 385 & 28.06 & 15.84 & 7.75 & $-$8.46 & 262 & 0.94 & 47.35 & 28.86 \\
61 & 0.086 & 103 & 385 & 28.28 & 15.98 & 7.89 & $-$8.59 & 259 & 0.94 & -- & -- \\
62 & 0.083 & 104 & 385 & 27.08 & 15.10 & 7.48 & $-$8.19 & 249 & 0.89 & -- & -- \\
63 & 0.080 & 105 & 385 & 25.64 & 14.16 & 6.77 & $-$7.37 & 239 & 0.84 & -- & -- \\
64 & 0.076 & 106 & 385 & 28.41 & 16.18 & 7.84 & $-$8.57 & 255 & 0.93 & -- & -- \\
65 & 0.073 & 107 & 386 & 29.27 & 16.66 & 7.81 & $-$8.55 & 271 & 0.97 & 25.62 & 14.21 \\
66 & 0.070 & 108 & 386 & 27.31 & 15.29 & 7.49 & $-$8.17 & 245 & 0.90 & -- & -- \\
67 & 0.068 & 109 & 386 & 27.22 & 15.15 & 7.43 & $-$8.11 & 245 & 0.90 & -- & -- \\
68 & 0.065 & 110 & 387 & 27.21 & 15.27 & 7.42 & $-$8.06 & 247 & 0.90 & -- & -- \\
69 & 0.062 & 111 & 387 & 27.89 & 15.81 & 7.72 & $-$8.45 & 241 & 0.92 & -- & -- \\
70 & 0.060 & 112 & 387 & 26.78 & 14.92 & 7.26 & $-$7.91 & 235 & 0.89 & 25.90 & 14.28 \\
71 & 0.057 & 113 & 387 & 26.98 & 15.03 & 7.38 & $-$7.99 & 255 & 0.91 & -- & -- \\
72 & 0.055 & 114 & 387 & 28.52 & 16.00 & 7.85 & $-$8.51 & 246 & 0.95 & -- & -- \\
73 & 0.053 & 115 & 387 & 25.51 & 14.08 & 6.86 & $-$7.45 & 226 & 0.84 & -- & -- \\
74 & 0.051 & 116 & 387 & 27.81 & 15.66 & 7.68 & $-$8.41 & 242 & 0.91 & -- & -- \\
75 & 0.050 & 117 & 387 & 24.92 & 13.57 & 6.57 & $-$7.10 & 239 & 0.83 & 68.64 & 45.86 \\
76 & 0.050 & 118 & 387 & 27.85 & 15.62 & 7.69 & $-$8.36 & 257 & 0.92 & -- & -- \\
77 & 0.050 & 119 & 387 & 28.85 & 16.40 & 8.07 & $-$8.81 & 248 & 0.94 & -- & -- \\
78 & 0.050 & 120 & 387 & 27.51 & 15.50 & 7.63 & $-$8.28 & 254 & 0.91 & -- & -- \\
79 & 0.050 & 121 & 387 & 26.71 & 14.90 & 7.34 & $-$8.03 & 239 & 0.87 & -- & -- \\
80 & 0.050 & 122 & 387 & 28.91 & 16.40 & 8.10 & $-$8.80 & 254 & 0.96 & 32.55 & 19.02 \\
\bottomrule
\end{tabular}
\caption[Escenario cruce2: episodios de entrenamiento (parte 2 de 3)]{Escenario cruce2: parámetros de entrada y métricas de salida por episodio (parte 2 de 3). 120 episodios de entrenamiento con semilla base 42; la fila fijo es el tiempo fijo con la misma semilla base. La política congelada se evalúa cada 5 episodios con una semilla fija ajena al entrenamiento y a la evaluación final.}
\label{tab:cruce2-episodios-2}
\end{table}

\begin{table}[H]
\setstretch{1}
\centering
\scriptsize
\setlength{\tabcolsep}{3pt}
\begin{tabular}{rrrrrrrrrrrr}
\toprule
\multicolumn{4}{c}{Entradas del episodio} & \multicolumn{6}{c}{Salidas del episodio (entrenamiento, $\varepsilon>0$)} & \multicolumn{2}{c}{Política congelada} \\
\cmidrule(lr){1-4}\cmidrule(lr){5-10}\cmidrule(lr){11-12}
Ep. & $\varepsilon$ & Semilla & Q ini. & Retraso & Espera & Cola & $r$/dec. & Cambios & Paradas & Retraso & Espera \\
 &  &  & (estados) & (s/veh) & (s/veh) & (veh) &  & de fase & (/veh) & (s/veh) & (s/veh) \\
\midrule
81 & 0.050 & 123 & 387 & 27.44 & 15.32 & 7.45 & $-$8.12 & 253 & 0.90 & -- & -- \\
82 & 0.050 & 124 & 387 & 28.37 & 16.04 & 7.81 & $-$8.52 & 251 & 0.95 & -- & -- \\
83 & 0.050 & 125 & 387 & 26.27 & 14.53 & 7.08 & $-$7.69 & 249 & 0.88 & -- & -- \\
84 & 0.050 & 126 & 387 & 27.55 & 15.58 & 7.61 & $-$8.34 & 233 & 0.90 & -- & -- \\
85 & 0.050 & 127 & 388 & 25.53 & 14.12 & 6.83 & $-$7.43 & 234 & 0.84 & 25.50 & 14.01 \\
86 & 0.050 & 128 & 388 & 25.88 & 14.40 & 6.94 & $-$7.57 & 233 & 0.85 & -- & -- \\
87 & 0.050 & 129 & 388 & 25.26 & 13.82 & 6.77 & $-$7.30 & 238 & 0.85 & -- & -- \\
88 & 0.050 & 130 & 388 & 27.21 & 15.07 & 7.41 & $-$8.02 & 257 & 0.90 & -- & -- \\
89 & 0.050 & 131 & 388 & 28.25 & 16.01 & 7.78 & $-$8.51 & 245 & 0.92 & -- & -- \\
90 & 0.050 & 132 & 389 & 26.26 & 14.46 & 7.03 & $-$7.64 & 253 & 0.88 & 31.05 & 17.48 \\
91 & 0.050 & 133 & 389 & 26.83 & 15.04 & 7.36 & $-$8.04 & 235 & 0.87 & -- & -- \\
92 & 0.050 & 134 & 390 & 25.59 & 14.07 & 6.86 & $-$7.42 & 235 & 0.86 & -- & -- \\
93 & 0.050 & 135 & 390 & 28.51 & 16.16 & 7.97 & $-$8.72 & 247 & 0.93 & -- & -- \\
94 & 0.050 & 136 & 390 & 27.80 & 15.61 & 7.60 & $-$8.26 & 253 & 0.93 & -- & -- \\
95 & 0.050 & 137 & 390 & 29.78 & 17.11 & 8.33 & $-$9.12 & 249 & 0.97 & 25.27 & 13.99 \\
96 & 0.050 & 138 & 390 & 26.39 & 14.77 & 7.17 & $-$7.82 & 236 & 0.86 & -- & -- \\
97 & 0.050 & 139 & 390 & 27.05 & 15.23 & 7.32 & $-$8.01 & 235 & 0.87 & -- & -- \\
98 & 0.050 & 140 & 390 & 26.72 & 14.85 & 7.23 & $-$7.84 & 249 & 0.89 & -- & -- \\
99 & 0.050 & 141 & 390 & 25.49 & 14.02 & 6.89 & $-$7.49 & 233 & 0.83 & -- & -- \\
100 & 0.050 & 142 & 390 & 26.94 & 15.05 & 7.39 & $-$8.08 & 245 & 0.90 & 25.28 & 13.99 \\
101 & 0.050 & 143 & 390 & 25.58 & 14.02 & 6.85 & $-$7.45 & 237 & 0.85 & -- & -- \\
102 & 0.050 & 144 & 390 & 25.46 & 14.03 & 6.83 & $-$7.43 & 233 & 0.85 & -- & -- \\
103 & 0.050 & 145 & 390 & 26.43 & 14.65 & 7.13 & $-$7.81 & 241 & 0.88 & -- & -- \\
104 & 0.050 & 146 & 390 & 27.00 & 15.13 & 7.35 & $-$8.03 & 245 & 0.89 & -- & -- \\
105 & 0.050 & 147 & 390 & 27.42 & 15.49 & 7.54 & $-$8.26 & 237 & 0.89 & 30.19 & 17.22 \\
106 & 0.050 & 148 & 390 & 25.72 & 14.21 & 6.92 & $-$7.52 & 229 & 0.84 & -- & -- \\
107 & 0.050 & 149 & 390 & 27.30 & 15.44 & 7.51 & $-$8.23 & 241 & 0.90 & -- & -- \\
108 & 0.050 & 150 & 390 & 26.21 & 14.57 & 7.13 & $-$7.78 & 233 & 0.86 & -- & -- \\
109 & 0.050 & 151 & 390 & 26.31 & 14.76 & 7.13 & $-$7.81 & 227 & 0.85 & -- & -- \\
110 & 0.050 & 152 & 390 & 25.45 & 14.10 & 6.83 & $-$7.49 & 224 & 0.82 & 29.60 & 16.72 \\
111 & 0.050 & 153 & 390 & 25.64 & 14.16 & 6.89 & $-$7.49 & 236 & 0.85 & -- & -- \\
112 & 0.050 & 154 & 390 & 27.74 & 15.57 & 7.61 & $-$8.29 & 241 & 0.92 & -- & -- \\
113 & 0.050 & 155 & 390 & 25.85 & 14.40 & 6.90 & $-$7.55 & 221 & 0.83 & -- & -- \\
114 & 0.050 & 156 & 390 & 27.46 & 15.46 & 7.48 & $-$8.19 & 233 & 0.89 & -- & -- \\
115 & 0.050 & 157 & 390 & 25.99 & 14.40 & 7.06 & $-$7.66 & 237 & 0.85 & 25.72 & 14.36 \\
116 & 0.050 & 158 & 390 & 25.60 & 14.12 & 6.87 & $-$7.48 & 237 & 0.84 & -- & -- \\
117 & 0.050 & 159 & 390 & 27.24 & 15.25 & 7.44 & $-$8.10 & 245 & 0.91 & -- & -- \\
118 & 0.050 & 160 & 390 & 26.53 & 14.74 & 7.19 & $-$7.80 & 237 & 0.87 & -- & -- \\
119 & 0.050 & 161 & 390 & 28.10 & 15.80 & 7.73 & $-$8.43 & 249 & 0.94 & -- & -- \\
120 & 0.050 & 162 & 390 & 26.53 & 14.81 & 7.19 & $-$7.84 & 238 & 0.87 & 179.00 & 134.62 \\
\bottomrule
\end{tabular}
\caption[Escenario cruce2: episodios de entrenamiento (parte 3 de 3)]{Escenario cruce2: parámetros de entrada y métricas de salida por episodio (parte 3 de 3). 120 episodios de entrenamiento con semilla base 42; la fila fijo es el tiempo fijo con la misma semilla base. La política congelada se evalúa cada 5 episodios con una semilla fija ajena al entrenamiento y a la evaluación final.}
\label{tab:cruce2-episodios-3}
\end{table}


% Generado por tablas_tex.py a partir de los CSV del proyecto. No editar a mano.
\begin{table}[H]
\setstretch{1}
\centering
\scriptsize
\setlength{\tabcolsep}{3pt}
\begin{tabular}{rrrrrrrrrrrr}
\toprule
\multicolumn{4}{c}{Entradas del episodio} & \multicolumn{6}{c}{Salidas del episodio (entrenamiento, $\varepsilon>0$)} & \multicolumn{2}{c}{Política congelada} \\
\cmidrule(lr){1-4}\cmidrule(lr){5-10}\cmidrule(lr){11-12}
Ep. & $\varepsilon$ & Semilla & Q ini. & Retraso & Espera & Cola & $r$/dec. & Cambios & Paradas & Retraso & Espera \\
 &  &  & (estados) & (s/veh) & (s/veh) & (veh) &  & de fase & (/veh) & (s/veh) & (s/veh) \\
\midrule
fijo & 0 & 42 & -- & 31.39 & 14.04 & 7.22 & $-$3.81 & 265 & 1.01 & -- & -- \\
\midrule
1 & 1.000 & 43 & 0 & 86.90 & 38.34 & 17.21 & $-$9.30 & 455 & 3.50 & -- & -- \\
2 & 0.900 & 44 & 133 & 61.99 & 26.70 & 12.51 & $-$6.64 & 450 & 2.60 & -- & -- \\
3 & 0.810 & 45 & 150 & 85.29 & 32.89 & 15.33 & $-$8.12 & 464 & 3.55 & -- & -- \\
4 & 0.729 & 46 & 165 & 54.30 & 21.40 & 10.18 & $-$5.27 & 457 & 2.38 & -- & -- \\
5 & 0.656 & 47 & 172 & 34.42 & 13.08 & 6.77 & $-$3.43 & 447 & 1.52 & 24.74 & 8.19 \\
6 & 0.590 & 48 & 176 & 47.34 & 18.83 & 9.33 & $-$4.86 & 455 & 2.09 & -- & -- \\
7 & 0.531 & 49 & 179 & 42.04 & 16.74 & 8.43 & $-$4.37 & 444 & 1.79 & -- & -- \\
8 & 0.478 & 50 & 183 & 37.35 & 14.26 & 7.28 & $-$3.72 & 450 & 1.68 & -- & -- \\
9 & 0.430 & 51 & 185 & 55.86 & 20.55 & 10.19 & $-$5.27 & 452 & 2.40 & -- & -- \\
10 & 0.387 & 52 & 188 & 40.19 & 14.72 & 7.48 & $-$3.82 & 456 & 1.78 & 30.09 & 10.90 \\
11 & 0.349 & 53 & 192 & 37.30 & 14.05 & 7.29 & $-$3.71 & 456 & 1.66 & -- & -- \\
12 & 0.314 & 54 & 193 & 26.84 & 9.69 & 5.15 & $-$2.57 & 442 & 1.15 & -- & -- \\
13 & 0.282 & 55 & 193 & 30.70 & 11.35 & 5.91 & $-$2.98 & 461 & 1.36 & -- & -- \\
14 & 0.254 & 56 & 194 & 30.45 & 11.55 & 6.02 & $-$3.06 & 466 & 1.33 & -- & -- \\
15 & 0.229 & 57 & 194 & 35.30 & 13.02 & 6.75 & $-$3.40 & 472 & 1.64 & 30.45 & 10.74 \\
16 & 0.206 & 58 & 195 & 31.78 & 11.81 & 6.18 & $-$3.10 & 450 & 1.39 & -- & -- \\
17 & 0.185 & 59 & 196 & 27.17 & 9.41 & 5.03 & $-$2.50 & 462 & 1.22 & -- & -- \\
18 & 0.167 & 60 & 196 & 29.16 & 10.79 & 5.70 & $-$2.85 & 465 & 1.28 & -- & -- \\
19 & 0.150 & 61 & 197 & 28.32 & 9.93 & 5.27 & $-$2.63 & 462 & 1.25 & -- & -- \\
20 & 0.135 & 62 & 197 & 29.06 & 10.06 & 5.30 & $-$2.64 & 472 & 1.31 & 23.89 & 7.91 \\
21 & 0.122 & 63 & 197 & 30.29 & 10.94 & 5.75 & $-$2.88 & 459 & 1.37 & -- & -- \\
22 & 0.109 & 64 & 197 & 27.97 & 10.87 & 5.75 & $-$2.93 & 420 & 1.16 & -- & -- \\
23 & 0.098 & 65 & 198 & 30.60 & 10.87 & 5.70 & $-$2.85 & 470 & 1.39 & -- & -- \\
24 & 0.089 & 66 & 201 & 37.55 & 13.27 & 6.80 & $-$3.45 & 464 & 1.67 & -- & -- \\
25 & 0.080 & 67 & 201 & 27.95 & 9.90 & 5.28 & $-$2.63 & 441 & 1.22 & 30.14 & 12.26 \\
26 & 0.072 & 68 & 201 & 27.77 & 10.02 & 5.29 & $-$2.64 & 452 & 1.21 & -- & -- \\
27 & 0.065 & 69 & 201 & 26.39 & 9.08 & 4.86 & $-$2.42 & 464 & 1.14 & -- & -- \\
28 & 0.058 & 70 & 201 & 30.55 & 10.55 & 5.56 & $-$2.77 & 462 & 1.42 & -- & -- \\
29 & 0.052 & 71 & 201 & 38.56 & 12.95 & 6.65 & $-$3.38 & 448 & 1.75 & -- & -- \\
30 & 0.050 & 72 & 201 & 42.91 & 14.94 & 7.62 & $-$3.87 & 448 & 1.96 & 27.97 & 9.69 \\
31 & 0.050 & 73 & 203 & 26.17 & 9.07 & 4.88 & $-$2.43 & 444 & 1.12 & -- & -- \\
32 & 0.050 & 74 & 203 & 32.18 & 10.78 & 5.71 & $-$2.84 & 458 & 1.45 & -- & -- \\
33 & 0.050 & 75 & 203 & 26.34 & 9.12 & 4.89 & $-$2.43 & 466 & 1.15 & -- & -- \\
34 & 0.050 & 76 & 203 & 26.81 & 9.13 & 4.90 & $-$2.43 & 463 & 1.18 & -- & -- \\
35 & 0.050 & 77 & 203 & 28.43 & 9.89 & 5.23 & $-$2.61 & 444 & 1.25 & 24.34 & 7.98 \\
36 & 0.050 & 78 & 203 & 26.78 & 9.43 & 5.01 & $-$2.49 & 451 & 1.17 & -- & -- \\
37 & 0.050 & 79 & 203 & 26.24 & 8.86 & 4.76 & $-$2.35 & 447 & 1.16 & -- & -- \\
38 & 0.050 & 80 & 203 & 27.02 & 9.29 & 4.96 & $-$2.45 & 455 & 1.20 & -- & -- \\
39 & 0.050 & 81 & 204 & 29.00 & 10.25 & 5.40 & $-$2.70 & 456 & 1.29 & -- & -- \\
40 & 0.050 & 82 & 204 & 29.29 & 10.60 & 5.59 & $-$2.81 & 448 & 1.27 & 30.14 & 10.81 \\
\bottomrule
\end{tabular}
\caption[Escenario corredor: episodios de entrenamiento]{Escenario corredor: parámetros de entrada y métricas de salida por episodio. 40 episodios de entrenamiento con semilla base 42; la fila fijo es el tiempo fijo con la misma semilla base. La política congelada se evalúa cada 5 episodios con una semilla fija ajena al entrenamiento y a la evaluación final.}
\label{tab:corredor-episodios}
\end{table}


% Generado por tablas_tex.py a partir de los CSV del proyecto. No editar a mano.
\begin{table}[H]
\setstretch{1}
\centering
\scriptsize
\setlength{\tabcolsep}{3pt}
\begin{tabular}{rrrrrrrrrrrr}
\toprule
\multicolumn{4}{c}{Entradas del episodio} & \multicolumn{6}{c}{Salidas del episodio (entrenamiento, $\varepsilon>0$)} & \multicolumn{2}{c}{Política congelada} \\
\cmidrule(lr){1-4}\cmidrule(lr){5-10}\cmidrule(lr){11-12}
Ep. & $\varepsilon$ & Semilla & Q ini. & Retraso & Espera & Cola & $r$/dec. & Cambios & Paradas & Retraso & Espera \\
 &  &  & (estados) & (s/veh) & (s/veh) & (veh) &  & de fase & (/veh) & (s/veh) & (s/veh) \\
\midrule
fijo & 0 & 42 & -- & 22.23 & 7.35 & 4.03 & $-$1.43 & 489 & 0.60 & -- & -- \\
\midrule
1 & 1.000 & 43 & 0 & 30.94 & 12.28 & 6.73 & $-$2.41 & 629 & 1.15 & -- & -- \\
2 & 0.950 & 44 & 88 & 30.71 & 12.15 & 6.70 & $-$2.40 & 649 & 1.16 & -- & -- \\
3 & 0.902 & 45 & 108 & 30.71 & 12.23 & 6.70 & $-$2.40 & 637 & 1.14 & -- & -- \\
4 & 0.857 & 46 & 124 & 29.33 & 11.19 & 6.16 & $-$2.19 & 662 & 1.12 & -- & -- \\
5 & 0.815 & 47 & 128 & 29.32 & 10.97 & 6.14 & $-$2.18 & 660 & 1.13 & 25.47 & 8.07 \\
6 & 0.774 & 48 & 132 & 27.36 & 9.71 & 5.31 & $-$1.87 & 683 & 1.06 & -- & -- \\
7 & 0.735 & 49 & 134 & 27.81 & 9.89 & 5.41 & $-$1.91 & 678 & 1.09 & -- & -- \\
8 & 0.698 & 50 & 136 & 28.97 & 10.89 & 5.04 & $-$1.77 & 684 & 1.12 & -- & -- \\
9 & 0.663 & 51 & 139 & 28.04 & 10.03 & 5.33 & $-$1.87 & 690 & 1.08 & -- & -- \\
10 & 0.630 & 52 & 139 & 27.96 & 10.06 & 5.47 & $-$1.92 & 683 & 1.08 & 25.32 & 8.08 \\
11 & 0.599 & 53 & 140 & 27.50 & 9.61 & 5.44 & $-$1.90 & 688 & 1.09 & -- & -- \\
12 & 0.569 & 54 & 140 & 26.65 & 8.96 & 4.82 & $-$1.68 & 724 & 1.04 & -- & -- \\
13 & 0.540 & 55 & 140 & 26.96 & 9.29 & 5.01 & $-$1.76 & 717 & 1.08 & -- & -- \\
14 & 0.513 & 56 & 140 & 26.11 & 8.76 & 4.74 & $-$1.66 & 706 & 1.03 & -- & -- \\
15 & 0.488 & 57 & 140 & 27.57 & 9.68 & 5.30 & $-$1.87 & 708 & 1.11 & 24.88 & 7.67 \\
16 & 0.463 & 58 & 141 & 27.36 & 9.45 & 5.19 & $-$1.80 & 720 & 1.12 & -- & -- \\
17 & 0.440 & 59 & 141 & 26.94 & 9.22 & 5.01 & $-$1.76 & 713 & 1.10 & -- & -- \\
18 & 0.418 & 60 & 142 & 25.62 & 8.29 & 4.48 & $-$1.57 & 714 & 1.02 & -- & -- \\
19 & 0.397 & 61 & 142 & 26.69 & 8.98 & 4.72 & $-$1.66 & 730 & 1.08 & -- & -- \\
20 & 0.377 & 62 & 142 & 25.58 & 8.38 & 4.47 & $-$1.56 & 737 & 1.04 & 24.78 & 7.93 \\
21 & 0.358 & 63 & 142 & 25.81 & 8.32 & 4.72 & $-$1.65 & 725 & 1.06 & -- & -- \\
22 & 0.341 & 64 & 142 & 24.97 & 7.60 & 4.17 & $-$1.44 & 775 & 1.03 & -- & -- \\
23 & 0.324 & 65 & 143 & 26.29 & 8.77 & 4.89 & $-$1.71 & 748 & 1.09 & -- & -- \\
24 & 0.307 & 66 & 143 & 25.21 & 7.90 & 4.30 & $-$1.49 & 751 & 1.02 & -- & -- \\
25 & 0.292 & 67 & 143 & 26.52 & 8.95 & 4.92 & $-$1.71 & 724 & 1.09 & 24.29 & 7.19 \\
26 & 0.277 & 68 & 144 & 25.63 & 8.22 & 4.51 & $-$1.57 & 728 & 1.03 & -- & -- \\
27 & 0.264 & 69 & 144 & 25.73 & 8.22 & 4.46 & $-$1.54 & 719 & 1.06 & -- & -- \\
28 & 0.250 & 70 & 144 & 25.24 & 8.17 & 4.47 & $-$1.55 & 732 & 1.05 & -- & -- \\
29 & 0.238 & 71 & 144 & 25.54 & 8.32 & 4.31 & $-$1.48 & 748 & 1.05 & -- & -- \\
30 & 0.226 & 72 & 144 & 25.91 & 8.69 & 4.47 & $-$1.54 & 736 & 1.08 & 24.27 & 7.31 \\
31 & 0.215 & 73 & 144 & 25.08 & 7.87 & 4.36 & $-$1.51 & 744 & 1.03 & -- & -- \\
32 & 0.204 & 74 & 144 & 24.21 & 7.60 & 4.15 & $-$1.43 & 737 & 0.97 & -- & -- \\
33 & 0.194 & 75 & 144 & 25.64 & 8.36 & 4.51 & $-$1.56 & 751 & 1.04 & -- & -- \\
34 & 0.184 & 76 & 145 & 25.08 & 8.02 & 4.44 & $-$1.55 & 739 & 1.03 & -- & -- \\
35 & 0.175 & 77 & 145 & 24.72 & 7.72 & 4.26 & $-$1.49 & 707 & 1.00 & 24.44 & 7.65 \\
36 & 0.166 & 78 & 145 & 24.47 & 7.58 & 4.28 & $-$1.49 & 722 & 1.00 & -- & -- \\
37 & 0.158 & 79 & 145 & 25.03 & 7.94 & 4.21 & $-$1.46 & 732 & 1.00 & -- & -- \\
38 & 0.150 & 80 & 145 & 24.36 & 7.39 & 4.24 & $-$1.47 & 742 & 0.97 & -- & -- \\
39 & 0.142 & 81 & 145 & 25.99 & 8.85 & 4.35 & $-$1.50 & 741 & 1.05 & -- & -- \\
40 & 0.135 & 82 & 145 & 24.41 & 7.54 & 4.17 & $-$1.45 & 733 & 0.97 & 23.91 & 7.29 \\
\bottomrule
\end{tabular}
\caption[Escenario red: episodios de entrenamiento (parte 1 de 3)]{Escenario red: parámetros de entrada y métricas de salida por episodio (parte 1 de 3). 100 episodios de entrenamiento con semilla base 42; la fila fijo es el tiempo fijo con la misma semilla base. La política congelada se evalúa cada 5 episodios con una semilla fija ajena al entrenamiento y a la evaluación final.}
\label{tab:red-episodios}
\end{table}

\begin{table}[H]
\setstretch{1}
\centering
\scriptsize
\setlength{\tabcolsep}{3pt}
\begin{tabular}{rrrrrrrrrrrr}
\toprule
\multicolumn{4}{c}{Entradas del episodio} & \multicolumn{6}{c}{Salidas del episodio (entrenamiento, $\varepsilon>0$)} & \multicolumn{2}{c}{Política congelada} \\
\cmidrule(lr){1-4}\cmidrule(lr){5-10}\cmidrule(lr){11-12}
Ep. & $\varepsilon$ & Semilla & Q ini. & Retraso & Espera & Cola & $r$/dec. & Cambios & Paradas & Retraso & Espera \\
 &  &  & (estados) & (s/veh) & (s/veh) & (veh) &  & de fase & (/veh) & (s/veh) & (s/veh) \\
\midrule
41 & 0.129 & 83 & 145 & 25.14 & 8.10 & 4.36 & $-$1.52 & 747 & 1.04 & -- & -- \\
42 & 0.122 & 84 & 145 & 24.93 & 7.91 & 4.33 & $-$1.50 & 766 & 0.99 & -- & -- \\
43 & 0.116 & 85 & 145 & 25.19 & 8.14 & 4.47 & $-$1.55 & 754 & 1.01 & -- & -- \\
44 & 0.110 & 86 & 145 & 24.91 & 7.84 & 4.43 & $-$1.53 & 766 & 1.01 & -- & -- \\
45 & 0.105 & 87 & 145 & 25.28 & 8.08 & 4.27 & $-$1.48 & 769 & 1.06 & 25.49 & 7.90 \\
46 & 0.099 & 88 & 145 & 25.13 & 8.05 & 4.16 & $-$1.45 & 762 & 1.03 & -- & -- \\
47 & 0.094 & 89 & 145 & 24.81 & 7.93 & 4.23 & $-$1.47 & 761 & 1.00 & -- & -- \\
48 & 0.090 & 90 & 145 & 24.59 & 7.54 & 4.27 & $-$1.47 & 768 & 1.00 & -- & -- \\
49 & 0.085 & 91 & 145 & 24.36 & 7.64 & 4.13 & $-$1.43 & 757 & 0.98 & -- & -- \\
50 & 0.081 & 92 & 145 & 24.85 & 7.80 & 4.35 & $-$1.50 & 746 & 1.03 & 24.00 & 7.03 \\
51 & 0.077 & 93 & 145 & 24.08 & 7.34 & 4.03 & $-$1.39 & 756 & 0.96 & -- & -- \\
52 & 0.073 & 94 & 145 & 25.65 & 8.49 & 4.20 & $-$1.46 & 764 & 1.04 & -- & -- \\
53 & 0.069 & 95 & 145 & 24.99 & 8.15 & 4.33 & $-$1.51 & 765 & 1.01 & -- & -- \\
54 & 0.066 & 96 & 145 & 25.27 & 7.93 & 4.36 & $-$1.50 & 753 & 1.04 & -- & -- \\
55 & 0.063 & 97 & 145 & 24.73 & 7.73 & 4.20 & $-$1.45 & 768 & 1.01 & 24.72 & 7.71 \\
56 & 0.060 & 98 & 145 & 25.12 & 8.16 & 4.28 & $-$1.49 & 752 & 1.00 & -- & -- \\
57 & 0.057 & 99 & 145 & 24.60 & 7.87 & 3.98 & $-$1.38 & 748 & 0.99 & -- & -- \\
58 & 0.054 & 100 & 145 & 25.32 & 7.78 & 4.32 & $-$1.50 & 769 & 1.08 & -- & -- \\
59 & 0.051 & 101 & 145 & 26.26 & 8.86 & 4.24 & $-$1.45 & 773 & 1.12 & -- & -- \\
60 & 0.050 & 102 & 145 & 25.00 & 7.89 & 4.38 & $-$1.52 & 757 & 1.03 & 24.45 & 7.69 \\
61 & 0.050 & 103 & 146 & 24.96 & 8.11 & 4.32 & $-$1.50 & 747 & 1.02 & -- & -- \\
62 & 0.050 & 104 & 146 & 23.75 & 7.29 & 4.14 & $-$1.42 & 738 & 0.94 & -- & -- \\
63 & 0.050 & 105 & 146 & 23.99 & 7.74 & 4.03 & $-$1.39 & 735 & 0.91 & -- & -- \\
64 & 0.050 & 106 & 146 & 24.52 & 7.82 & 4.34 & $-$1.52 & 749 & 0.97 & -- & -- \\
65 & 0.050 & 107 & 146 & 25.31 & 8.04 & 4.36 & $-$1.50 & 761 & 1.03 & 24.36 & 7.51 \\
66 & 0.050 & 108 & 146 & 25.10 & 8.07 & 4.32 & $-$1.50 & 747 & 1.01 & -- & -- \\
67 & 0.050 & 109 & 146 & 24.45 & 7.52 & 4.29 & $-$1.49 & 751 & 0.99 & -- & -- \\
68 & 0.050 & 110 & 146 & 24.80 & 7.91 & 4.25 & $-$1.47 & 736 & 1.00 & -- & -- \\
69 & 0.050 & 111 & 146 & 23.89 & 7.29 & 4.03 & $-$1.40 & 749 & 0.94 & -- & -- \\
70 & 0.050 & 112 & 146 & 24.45 & 7.56 & 4.22 & $-$1.46 & 748 & 0.97 & 24.23 & 7.40 \\
71 & 0.050 & 113 & 148 & 25.40 & 8.30 & 4.34 & $-$1.49 & 751 & 1.06 & -- & -- \\
72 & 0.050 & 114 & 148 & 24.79 & 7.68 & 4.24 & $-$1.47 & 770 & 1.01 & -- & -- \\
73 & 0.050 & 115 & 148 & 24.20 & 7.58 & 4.25 & $-$1.47 & 764 & 1.00 & -- & -- \\
74 & 0.050 & 116 & 148 & 24.74 & 7.52 & 4.32 & $-$1.50 & 758 & 1.01 & -- & -- \\
75 & 0.050 & 117 & 148 & 24.86 & 7.74 & 4.25 & $-$1.47 & 750 & 1.02 & 23.85 & 7.15 \\
76 & 0.050 & 118 & 148 & 24.30 & 7.49 & 4.02 & $-$1.40 & 759 & 0.94 & -- & -- \\
77 & 0.050 & 119 & 148 & 24.24 & 7.45 & 4.09 & $-$1.42 & 751 & 0.98 & -- & -- \\
78 & 0.050 & 120 & 148 & 24.98 & 8.09 & 4.16 & $-$1.44 & 756 & 1.00 & -- & -- \\
79 & 0.050 & 121 & 148 & 23.67 & 7.29 & 4.00 & $-$1.39 & 759 & 0.94 & -- & -- \\
80 & 0.050 & 122 & 148 & 24.52 & 7.57 & 4.26 & $-$1.47 & 747 & 0.97 & 23.77 & 7.03 \\
\bottomrule
\end{tabular}
\caption[Escenario red: episodios de entrenamiento (parte 2 de 3)]{Escenario red: parámetros de entrada y métricas de salida por episodio (parte 2 de 3). 100 episodios de entrenamiento con semilla base 42; la fila fijo es el tiempo fijo con la misma semilla base. La política congelada se evalúa cada 5 episodios con una semilla fija ajena al entrenamiento y a la evaluación final.}
\label{tab:red-episodios-2}
\end{table}

\begin{table}[H]
\setstretch{1}
\centering
\scriptsize
\setlength{\tabcolsep}{3pt}
\begin{tabular}{rrrrrrrrrrrr}
\toprule
\multicolumn{4}{c}{Entradas del episodio} & \multicolumn{6}{c}{Salidas del episodio (entrenamiento, $\varepsilon>0$)} & \multicolumn{2}{c}{Política congelada} \\
\cmidrule(lr){1-4}\cmidrule(lr){5-10}\cmidrule(lr){11-12}
Ep. & $\varepsilon$ & Semilla & Q ini. & Retraso & Espera & Cola & $r$/dec. & Cambios & Paradas & Retraso & Espera \\
 &  &  & (estados) & (s/veh) & (s/veh) & (veh) &  & de fase & (/veh) & (s/veh) & (s/veh) \\
\midrule
81 & 0.050 & 123 & 148 & 24.61 & 7.71 & 4.26 & $-$1.47 & 748 & 1.00 & -- & -- \\
82 & 0.050 & 124 & 148 & 24.26 & 7.55 & 4.13 & $-$1.43 & 769 & 0.98 & -- & -- \\
83 & 0.050 & 125 & 148 & 24.80 & 7.83 & 4.39 & $-$1.54 & 772 & 1.02 & -- & -- \\
84 & 0.050 & 126 & 148 & 24.00 & 7.44 & 4.03 & $-$1.39 & 737 & 0.95 & -- & -- \\
85 & 0.050 & 127 & 148 & 24.59 & 7.73 & 4.25 & $-$1.47 & 748 & 1.00 & 24.46 & 7.56 \\
86 & 0.050 & 128 & 148 & 24.89 & 7.97 & 4.26 & $-$1.48 & 740 & 0.97 & -- & -- \\
87 & 0.050 & 129 & 148 & 24.71 & 7.81 & 4.29 & $-$1.49 & 738 & 0.99 & -- & -- \\
88 & 0.050 & 130 & 148 & 23.99 & 7.22 & 4.15 & $-$1.45 & 731 & 0.96 & -- & -- \\
89 & 0.050 & 131 & 149 & 24.05 & 7.33 & 4.09 & $-$1.42 & 758 & 0.95 & -- & -- \\
90 & 0.050 & 132 & 149 & 24.63 & 7.69 & 4.22 & $-$1.47 & 766 & 0.99 & 24.38 & 7.56 \\
91 & 0.050 & 133 & 149 & 24.64 & 7.76 & 4.36 & $-$1.52 & 755 & 0.99 & -- & -- \\
92 & 0.050 & 134 & 149 & 24.70 & 7.62 & 4.24 & $-$1.47 & 748 & 1.03 & -- & -- \\
93 & 0.050 & 135 & 149 & 25.71 & 8.58 & 4.47 & $-$1.53 & 758 & 1.04 & -- & -- \\
94 & 0.050 & 136 & 149 & 24.56 & 7.89 & 4.15 & $-$1.44 & 755 & 0.96 & -- & -- \\
95 & 0.050 & 137 & 149 & 24.30 & 7.65 & 4.11 & $-$1.43 & 758 & 0.95 & 24.85 & 7.73 \\
96 & 0.050 & 138 & 149 & 25.11 & 7.99 & 4.37 & $-$1.52 & 775 & 1.03 & -- & -- \\
97 & 0.050 & 139 & 149 & 24.62 & 7.69 & 4.16 & $-$1.45 & 775 & 0.98 & -- & -- \\
98 & 0.050 & 140 & 149 & 25.48 & 8.22 & 4.15 & $-$1.43 & 774 & 1.02 & -- & -- \\
99 & 0.050 & 141 & 149 & 24.94 & 7.98 & 4.45 & $-$1.54 & 765 & 1.01 & -- & -- \\
100 & 0.050 & 142 & 149 & 25.36 & 8.16 & 4.58 & $-$1.59 & 761 & 1.05 & 24.84 & 7.58 \\
\bottomrule
\end{tabular}
\caption[Escenario red: episodios de entrenamiento (parte 3 de 3)]{Escenario red: parámetros de entrada y métricas de salida por episodio (parte 3 de 3). 100 episodios de entrenamiento con semilla base 42; la fila fijo es el tiempo fijo con la misma semilla base. La política congelada se evalúa cada 5 episodios con una semilla fija ajena al entrenamiento y a la evaluación final.}
\label{tab:red-episodios-3}
\end{table}



\section*{Anexo B. Código fuente}

Scripts del prototipo tal como se ejecutaron para producir los
resultados del Capítulo~\ref{cap:resultados} (SUMO 1.25.0, Python 3.14).
El código, sus comentarios y sus salidas están en español y sin tildes
en los identificadores. Se omite \texttt{graficar.py}, que solo dibuja
las figuras a partir de los mismos archivos; está en el repositorio del
proyecto.

\begin{spacing}{1}

\subsection*{rl\_semaforo.py: agente único (escenarios cruce y cruce2)}
{\scriptsize \verbatiminput{codigo/rl_semaforo.py}}

\subsection*{rl\_corredor.py: un agente independiente por semáforo (corredor y red)}
{\scriptsize \verbatiminput{codigo/rl_corredor.py}}

\subsection*{comun.py: métricas del tripinfo, programa actuado y CSV}
{\scriptsize \verbatiminput{codigo/comun.py}}

\subsection*{evaluar.py: evaluación por semillas y estadística}
{\scriptsize \verbatiminput{codigo/evaluar.py}}

\subsection*{barrer\_fijo.py: barrido del programa de tiempo fijo}
{\scriptsize \verbatiminput{codigo/barrer_fijo.py}}

\subsection*{tablas\_tex.py: tablas del capítulo de resultados}
{\scriptsize \verbatiminput{codigo/tablas_tex.py}}

\end{spacing}
